\begin{abstract}
Complex-valued signals, such as Magnetic Resonance Imaging (MRI) and audio spectrograms, are almost universally modelled as flat two-channel Euclidean data. For nonzero values, the amplitude--phase chart $z \mapsto (|z|, z/|z|)$ identifies the signal domain with the cylinder $(0, \infty) \times S^1$, on which we deliberately replace the inherited metric $dA^2 + A^2 d\theta^2$ by the decoupled product metric $dA^2 + d\theta^2$. In this empirical study, we measure what that substitution costs and what it buys. Exact analytical bridges computed across three distinct sources of complex data---synthetic copula fields, coil-combined fastMRI knee acquisitions, and LibriSpeech spectrograms---demonstrate that Cartesian paths induce a heavy-tailed distribution of angular velocity with a Pareto tail index of $\approx 1$. Under independent coupling, $43$--$49\%$ of the total signal energy falls on paths that turn faster than $\pi$ rad per unit time, a rotational speed no cylindrical path ever reaches. To resolve this singularity, we formulate Cylindrical Flow Matching (CyFM), which strictly bounds the angular regression target, and couple noise and data via exact minibatch Optimal Transport computed jointly over whole fields in the cylindrical metric. On synthetic fields, this joint coupling reduces few-step generation error by $3$--$60\%$. Against the strongest Cartesian baseline, CyFM achieves lower generative error at every integration step up to $k = 8$ across all evaluated spatial resolutions with separated seeds, as well as on real speech spectrograms, while exhibiting no significant difference at asymptotic convergence ($k = 100$). Finally, a prior-only baseline underscores the severe failure of flat models in few-step regimes: on synthetic fields, a single Cartesian Euler step performs worse than the unintegrated noise prior ($0.376$ vs.\ $0.150$).
\end{abstract}
